\section{Marco Teórico}
\label{sec:marco}

\subsection{Grafos y análisis espectral}

Un grafo $G = (V, E, w)$ es una estructura matemática compuesta por un
conjunto de nodos $V$, un conjunto de aristas $E$ y una función de peso
$w : E \to \mathbb{R}$. En este trabajo los nodos representan tokens de
criptoactivos y los pesos reflejan el grado de correlación entre sus precios.

La \emph{Laplaciana normalizada} $\mathcal{L} = D^{-1/2}(D-A)D^{-1/2}$
tiene eigenvalores acotados en $[0,2]$ independientemente de la
distribución de pesos, lo que hace que los valores del IPR sean
comparables entre ventanas de tiempo distintas --- propiedad esencial
para construir una serie temporal coherente del detector. La Laplaciana
estándar $D - A$ no tiene esta propiedad: su espectro escala con el
grado medio de la red, de modo que ventanas con distinta densidad de
correlaciones producen rangos de eigenvalores incomparables. La
matriz de transición de caminata aleatoria $D^{-1}A$ tampoco es simétrica,
lo que impide garantizar que sus eigenvectores formen una base ortonormal y
complica el cálculo del IPR~\cite{newman2003}.

La normalización por $D^{-1/2}$ sí amplifica la contribución de nodos
periféricos (tokens con pocas aristas fuertes), pero este efecto es
deseable en nuestro contexto: un token aislado que de repente forma
correlaciones fuertes con un grupo es precisamente la señal que queremos
detectar. El ruido que puede introducirse en nodos muy aislados queda
controlado por el filtrado espectral previo con RMT, que elimina
correlaciones espurias antes de construir el grafo.

El \emph{Inverse Participation Ratio} (IPR) de un eigenvector mide qué
tan concentrada está su energía sobre los nodos. Un eigenvector
\emph{extendido} produce $\mathrm{IPR} \approx N^{-1}$; uno
\emph{localizado} produce $\mathrm{IPR} \gg N^{-1}$. Este concepto
proviene de la física del estado sólido~\cite{anderson1958absence}, donde
se usa para estudiar la transición entre estados conductores (extendidos)
y aislantes (localizados) en materiales desordenados.

\subsection{Esparsificación de grafos}

Un grafo completo sobre $N$ activos tiene $N(N-1)/2$ aristas. Una red tan
densa uniformiza los eigenvectores de la Laplaciana, llevando el IPR de
todos los modos cerca de $N^{-1}$ y borrando cualquier señal de
localización.

El \emph{Árbol de Expansión Mínima} (MST) conecta todos los nodos con
exactamente $N-1$ aristas de mínimo costo total, preservando las
correlaciones más fuertes~\cite{mantegna1999}. Tumminello \textit{et
al.}~\cite{tumminello2005} demostraron que estructuras de filtrado
similares retienen información estadísticamente significativa en mercados
financieros. El MST se usa como esqueleto base al que se añaden aristas
adicionales calibradas por modularidad, equilibrando parsimonia
estructural con riqueza topológica.

\subsection{Mercados financieros como redes de correlación}

Los retornos logarítmicos son la representación estándar del cambio de
precio de un activo porque estabilizan la varianza y hacen comparables
activos de distinto precio. La correlación de Pearson mide dependencia
lineal y es el punto de partida estándar en la literatura de redes
financieras~\cite{mantegna1999,plerou1999}.

Sin embargo, los retornos de criptoactivos exhiben colas pesadas extremas
y agrupamiento de volatilidad (\emph{volatility clustering}), fenómenos
que pueden inflar artificialmente la correlación de Pearson durante
choques exógenos~\cite{xu2019anatomy}. Medidas de dependencia basadas en
rangos, como Spearman o Kendall, son más robustas a valores atípicos
porque solo dependen del orden relativo de los retornos y no de su
magnitud. No las adoptamos como métrica principal porque la RMT
---que es el mecanismo de filtrado central del pipeline--- está
desarrollada teóricamente para matrices de correlación de Pearson; no
existe un equivalente de Marchenko-Pastur para matrices de correlación
de Spearman con las mismas propiedades analíticas. No obstante, la
sensibilidad a este supuesto se evaluará empíricamente comparando los
resultados del detector bajo Pearson y bajo Spearman como análisis de
robustez.

Mantegna~\cite{mantegna1999} propuso transformar las correlaciones en
distancias mediante $d_{ij} = \sqrt{2(1 - C_{ij})}$, lo que define una
métrica euclidiana sobre el espacio de activos y permite construir el
MST con semántica geométrica clara.

En mercados de criptoactivos existe un fenómeno dominante: casi todos los
tokens se mueven junto con Bitcoin. Esta covariación se manifiesta como
un eigenvalor dominante de la matriz de correlación que corresponde al
\emph{modo de mercado global}~\cite{plerou1999}. Si no se elimina,
enmascara cualquier estructura de comunidad más fina.

\subsection{Teoría de Matrices Aleatorias y estacionariedad}

La RMT predice que, bajo correlaciones puramente aleatorias, los
eigenvalores de una matriz de correlación $N \times N$ construida sobre
series de longitud $T$ caen dentro del intervalo de Marchenko-Pastur
$[\lambda_-, \lambda_+]$, que depende únicamente del cociente $Q =
N/T$~\cite{plerou1999}. Los eigenvalores fuera de este intervalo
contienen información genuina sobre la estructura del sistema.

La condición $T \gg N$ (equivalentemente, $Q \ll 1$) es necesaria para
que la predicción teórica sea precisa. Con velas de 1~minuto, la ventana
de calibración de $T_{\text{cal}} = 720~\text{h}$ corresponde a
$720 \times 60 = 43\,200$ observaciones. Para $N = 50$ tokens activos,
el cociente es $Q = 50/43\,200 \approx 0.0012$, valor que satisface
ampliamente la condición $Q \ll 1$ y justifica la aplicación de la
distribución de Marchenko-Pastur. En cambio, la ventana de detección de
$T_{\text{det}} = 48~\text{h}$ produce solo $2\,880$ observaciones,
con $Q \approx 0.017$, donde el estimador de Pearson tiene sesgo
significativo. La solución adoptada separa la estimación de los límites
(ventana larga, RMT válida) de su aplicación sobre ventanas cortas.

Esta transferencia introduce un supuesto de \emph{estacionariedad local}:
que la estructura de ruido estimada en los 30 días previos es
representativa de las siguientes 48~h. El mercado de criptoactivos sufre
cambios de régimen bruscos que pueden violar este supuesto. Sin embargo,
el umbral dinámico basado en el Configuration Model (véase
Sección~\ref{sec:solucion}) mitiga parcialmente este problema: al
recalibrarse en cada ventana de detección usando la topología local de
$G_t$, el umbral absorbe cambios en el nivel de correlación agregado sin
requerir estacionariedad global. La validez de la transferencia se
verifica empíricamente mediante un análisis de sensibilidad descrito en
la sección de solución propuesta.

\subsection{Modelos nulos en redes complejas}

Detectar una anomalía requiere definir primero qué constituye el
comportamiento normal. En redes complejas el \emph{modelo nulo} estándar
es el \emph{Configuration Model}~\cite{newman2003}: dado un grafo
observado $G$ con secuencia de grados $\{d_i\}_{i=1}^N$, el
Configuration Model genera una familia de grafos aleatorios que preservan
exactamente esa secuencia de grados pero con las aristas reasignadas
uniformemente al azar.

Esta construcción tiene una propiedad crucial: al preservar la secuencia
de grados, controla por la densidad local de la red. Un IPR alto no
puede atribuirse simplemente a que la red sea densa o esté concentrada
en pocos hubs --- esas características ya están capturadas por el modelo
nulo. Lo que el umbral detecta es un exceso de localización que no puede
explicarse por la topología de grado sola.

El umbral de detección se define como el percentil~99.7 empírico de la
distribución de $\mathrm{IPR}_{p90}$ sobre las $M$ redes nulas
generadas. Este criterio no paramétrico es preferible a $\mu_{\text{null}}
+ 3\,\sigma_{\text{null}}$ porque el estadístico $\mathrm{IPR}_{p90}$
---siendo un estadístico de orden sobre una métrica de red--- puede
seguir una distribución asimétrica bajo el Configuration Model. Asumir
normalidad subestimaría el umbral cuando la distribución nula tiene cola
derecha pesada, incrementando los falsos positivos.

Una limitación del Configuration Model puro es que destruye el
coeficiente de clustering local al reasignar aristas al azar. Las redes
financieras tienen alta transitividad natural: si el token $A$ correlaciona
fuertemente con $B$ y $B$ con $C$, es probable que $A$ correlacione con
$C$, formando triángulos. Un modelo nulo que no captura esta propiedad
tiene un $\mathrm{IPR}_{p90}$ nulo sistemáticamente menor que el de la
red real en condiciones no manipuladas, haciendo que el umbral del
99.7\% sea demasiado permisivo e incrementando los falsos positivos.

Para cuantificar este efecto se comparará el Configuration Model puro
con un modelo nulo de \emph{rewiring local}: dado el grafo observado
$G_t$, se intercambian repetidamente los extremos de pares de aristas
elegidos al azar preservando la secuencia de grados y el coeficiente de
clustering local dentro de una tolerancia del $5\%$~\cite{maslov2002}.
Al preservar la transitividad de la red, este modelo nulo produce un
umbral más conservador y reduce los falsos positivos atribuibles a la
estructura natural de triángulos del mercado. Los resultados se reportan
bajo ambos modelos nulos.

\subsection{Limitaciones del marco: coordinación artificial frente a
choques legítimos}

La hipótesis central de este trabajo es que la fase de acumulación de un
\emph{pump \& dump} produce localización espectral anómala en la red de
correlaciones. Sin embargo, no toda localización espectral implica
manipulación. Un choque fundamental legítimo y localizado ---como una
noticia regulatoria favorable exclusivamente para el sector DeFi--- haría
que ese grupo de tokens formara una comunidad densa y transitoria, elevando
$\mathrm{IPR}_{p90}$ de la misma forma que la manipulación.

El marco teórico no puede distinguir ambos casos usando solo información
espectral. La diferenciación depende de dos elementos que están fuera del
alcance de este trabajo: (i)~información externa sobre eventos de mercado
que permita etiquetar períodos de choque legítimo, y (ii)~patrones
temporales (duración, reversión brusca) que son distintos entre un
\emph{pump \& dump} y un choque fundamental. Esta es una limitación
conocida del enfoque y no invalida el detector; implica que en la
evaluación empírica se reportarán los falsos positivos coincidentes con
eventos externos identificables, separándolos de los errores atribuibles
al modelo.
